% D.80. Root, FGH, and PBC complete independent proof/source/certificate audits PASS.
\subsection{An exact compact-kernel family with arbitrary infinite ramification}

\begin{statement}{D.80}{Proved}{Unrestricted independent unit pairs and actual point examples}
Let \(a,b\in H^1(G_{\mathbb Q},\mathbb Z_5(1))\) be arbitrary
completed Kummer classes. Assume that their valuations at five
are zero, their reductions modulo five are linearly independent,
and their local unit-logarithm vector is primitive:
\[
 a|_{G_{\mathbb Q_5}}=\alpha\kappa(6),\qquad
 b|_{G_{\mathbb Q_5}}=\beta\kappa(6),\qquad
                  (\alpha,\beta)\notin5\mathbb Z_5^2.
                                                               \tag{D.80.1}
\]
No finite ramification set for \(a,b\) is assumed. Define
\[
 d=\max\{v_5(\alpha),v_5(\beta),v_5(\alpha-\beta)\},
                    \qquad v_5(0)=\infty.
                                                               \tag{D.80.2}
\]
Then we have the full compact-kernel classification
\[
 \mathcal S_s=
 \begin{cases}\mathbb Z/5,&d=0,\\0,&d>0,\end{cases}
 \qquad\mathcal T_s=0,\qquad
 \mathcal K_s\simeq
 \begin{cases}
   \mathbb Z/5,&d=0,\\
   \mathbb Z_5^d,&1\le d<\infty,\\
   \displaystyle\prod_{\mathbb N}\mathbb Z_5,&d=\infty.
 \end{cases}                                                   \tag{D.80.3}
\]
These are compact topological groups in D.63's separated
second-cohomology sequence. The general nonseparated kernel,
arbitrary infinitely ramified first-cohomology ranks, and
section existence are not determined by this assertion.

There are actual rational-point fibres with compact kernel
\(\mathbb Z_5^d\) for every integer \(d\ge1\). Explicitly choose
the integer \(t_d\) by
\[
 \begin{split}
 t_d&=5^{d+1}-1+5^{d+2}k_d,\qquad 0\le k_d<36,\\
 k_d&\equiv(23-5^{d+1})(5^{d+2})^{-1}\pmod{36}.
 \end{split}                                                   \tag{D.80.4}
\]
The actual point \(t_d\) on the tripod has pair
\((\kappa(t_d),\kappa(1-t_d))\) satisfying the hypotheses
and depth \(d\). Its nonempty metabelian section fibre is
a torsor under \(\prod_{\mathbb N}\mathbb Z_5\), by D.74.
The point \(t=-2\) supplies the order-five case \(d=0\).
\end{statement}

\begin{proof}
\paragraph{The actual image from finitely many residual inertia vectors.}
The pro-five completed multiplicative group of \(\mathbb Q\)
has the valuation description
\[
 H^1(G_{\mathbb Q},\mathbb Z_5(1))
   =\{(x_q)_q\in\prod_q\mathbb Z_5:
                     x_q\longrightarrow0\text{ five-adically}\}.
                                                               \tag{D.80.5}
\]
Indeed \(\mathbb Q^\times=\{\pm1\}\oplus\bigoplus_q\mathbb Z\)
and the sign disappears in the pro-five completion. At each
finite coefficient level its valuation support is finite.
The local Kummer valuation of a class is its indicated
coordinate, so (D.80.5) applies to actual inertia images.

Modulo five, the vectors \((a_q,b_q)\) for \(q\ne5\)
span \(\mathbb F_5^2\). Otherwise a nonzero linear relation
on that span, together with the zero coordinate at five,
would contradict the assumed independence of \(a,b\).
Only finitely many such vectors are nonzero modulo five.
Choose two whose determinant is a unit in \(\mathbb Z_5\).
Their actual inertia images lie in the pure translation
kernel of the untwisted operator image \(I_0\), and
generate the whole translation group \(H=\mathbb Z_5^2\).
Thus D.78's faithful compact-image proof gives
\[
                         I_s=H\rtimes I_0.             \tag{D.80.6}
\]
This does not require bounded global ramification.

The local logarithm in (D.80.1) is continuous on the completed
group. Explicitly, for \(x_5=0\) its coefficient is
\[
      \sum_{q\ne5}x_q\frac{\log(q^4)}{4\log6}
                          \quad\hbox{in }\mathbb Z_5.
                                                               \tag{D.80.7}
\]
The summands tend uniformly to zero because the logarithm
ratios are integral and \(x_q\to0\). Thus the local unit
description does not involve a conditionally convergent
choice of ordering. After an integral cocycle gauge,
the exact local image is
\[
 I_5=\mathbb Z_5(\alpha,\beta)\rtimes I_0,              \tag{D.80.8}
\]
by the same local surjectivity and \(\kappa(6)\)-saturation
argument of D.78. The chosen local line is primitive.

The coinduction argument, the finite-image Chebotarev
injection of \(\mathcal S_s^*\), and the complete local
restriction-kernel calculation of D.79 now apply to
(D.80.6)--(D.80.8). None uses the global support after
these actual images have been established. In particular
\[
 \mathcal S_s^*=0\quad(d>0),\qquad
 \mathcal S_s^*\subseteq\mathbb F_5\quad(d=0).           \tag{D.80.9}
\]
For \(d=0\), the possible generator is the cocycle
\(c=48h\bmod5\) with constant Cartier coefficients,
as in D.78. It remains to prove that this class occurs
at one common finite coefficient level despite possible
infinite ramification.

\paragraph{Local vanishing on discrete Cartier coefficients.}
Fix \(q\ne5\) at which the abelian inertia vector is
nonzero. Put \(A=A_s\), \(W_r=A/\mathfrak m^r A\),
and \(D=A^D=\varinjlim W_r^D\), where
\(\mathfrak m=(5,u,v)\). D.43 proves
\(H^1(G_{\mathbb Q_q},A)=0\) with compact integral
coefficients, for every such inertia vector, whether
or not a metabelian section exists.

Here is the exact passage to the discrete dual; it is
not an unsupported formal duality assertion. Each
\(H^0(G_{\mathbb Q_q},W_r)\) and
\(H^1(G_{\mathbb Q_q},W_r)\) is finite. The compact
coefficient exact sequence and the Mittag--Leffler
property of the first tower therefore give
\[
 \varprojlim_r H^1(G_{\mathbb Q_q},W_r)
                    =H^1(G_{\mathbb Q_q},A)=0.          \tag{D.80.10}
\]
Finite local Tate duality
\cite[I, Corollary 2.3]{DE_MilneADT} identifies
\(H^1(G_{\mathbb Q_q},W_r^D)\) with the finite
Pontryagin dual of \(H^1(G_{\mathbb Q_q},W_r)\),
naturally for coefficient transitions. Pontryagin
duality between finite inverse and direct limits, and
continuity of cohomology on discrete filtered unions,
then give
\[
 \begin{split}
 H^1(G_{\mathbb Q_q},D)
  &=\varinjlim_r H^1(G_{\mathbb Q_q},W_r^D)\\
  &=\left(\varprojlim_r H^1(G_{\mathbb Q_q},W_r)\right)^*
   =0.
 \end{split}                                             \tag{D.80.11}
\]
In particular any class represented at one finite level
becomes a coboundary at some larger finite level.
No uniformity in \(q\) is concluded from this step alone.

\paragraph{Only finitely many places need those larger levels.}
Use the one fixed module of D.78,
\[
          M=(A_s/(5,(u,v)^4)A_s)^D\subset W_4^D.
                                                               \tag{D.80.12}
\]
Let
\[
 E=\{q\ne5:(a_q,b_q)\notin5\mathbb Z_5^2\}.
                                                               \tag{D.80.13}
\]
This set is finite by (D.80.5). For every \(q\notin E\cup\{5\}\),
the module \(M\) is unramified even if \(A_s\) is ramified:
its inertia multiplier satisfies
\[
 t_1^{a_q}t_2^{b_q}\equiv1\pmod{5,(u,v)^4},             \tag{D.80.14}
\]
because \((1+u)^{5x}=(1+u^5)^x\) modulo five, and similarly
for \(v\). The original untwisted action and \(c\) are
unramified at every \(q\ne5\). Hence D.78's universal
constant, linear and cubic Frobenius primitives trivialize
the local class of \(c\) in \(M\) at every prime outside
this finite exceptional set. The identities apply to the
actual Frobenius values of \(a,b,c,\chi\), without a
restriction on their ramification at higher precision.

At each of the finitely many primes in \(E\), the inertia
vector is nonzero. Equation (D.80.11) kills the local
class at some finite level \(W_{r_q}^D\). At five,
the \(d=0\) unit-slope primitive already lies in the
three-jet module \(M\), by D.78. Real cohomology is
zero at this odd prime. Choose
\[
                         R\ge\max(4,\max_{q\in E}r_q).
                                                               \tag{D.80.15}
\]
The primitives originally lying in \(M\) remain primitives
in \(W_R^D\), even where the larger module is ramified.
Thus the single global class \([c]\) is everywhere locally
trivial in this one finite module \(W_R^D\). Its nonzero
image in \(H^1(G_{\mathbb Q},D)\) follows from the actual
coinduction isomorphism and primitive raw Soul\'e restriction,
exactly as in D.78. Therefore \(\mathcal S_s^*=\mathbb F_5\)
when \(d=0\). This proves the first part of (D.80.3)
without commuting a direct limit with an infinite product.

Independence modulo five makes \(a,b,a-b\) all globally
primitive. D.77 therefore gives \(\mathcal T_s=0\), and
its exact local-depth calculation and continuous product
decomposition give the remaining parts of (D.80.3).

\paragraph{Actual ordinary fibres of every finite free rank.}
The inverse in (D.80.4) exists since five is prime to \(36\).
The resulting integer has
\[
 t_d\equiv22\pmod{36},\qquad
 t_d\equiv-1+5^{d+1}\pmod{5^{d+2}}.                    \tag{D.80.16}
\]
It is different from zero and one. The first congruence
gives \(v_2(t_d)=1\), \(v_3(1-t_d)=1\), while
\(v_2(1-t_d)=v_3(t_d)=0\). Thus the two global Kummer
classes are linearly independent modulo five, detected
by actual inertia at two and three. Both are units at five.

For \(d\ge1\), the second congruence gives
\[
 v_5(t_d^4-1)=d+1,\qquad 1-t_d\equiv2\pmod{25}.
\]
The five-adic logarithm preserves the positive valuation
of a principal unit minus one. Since \(v_5(\log6)=1\),
the local coefficient \(\alpha\) has valuation exactly
\(d\), and \(\beta\equiv2\pmod5\). Their difference
is a unit. The pair therefore has exactly the depth
claimed in (D.80.2), proving \(\mathcal K_s=\mathbb Z_5^d\).
The ordinary rational point supplies an actual full
section; its metabelian fibre is nonempty and the
finite-support statement D.74 applies. This proves
the final assertion, without asserting point-theoreticity
of every class in that much larger fibre.
\end{proof}
